\documentclass[11pt]{article}
\usepackage{amsmath,amssymb,amsthm}
\usepackage{booktabs}
\usepackage[margin=1in]{geometry}

\newtheorem{theorem}{Theorem}
\newtheorem{lemma}[theorem]{Lemma}
\newtheorem{corollary}[theorem]{Corollary}

\title{Problem 759: A Squared Recurrence}
\author{Project Euler Solutions}
\date{}

\begin{document}
\maketitle


\section{Problem Statement}
$f(1)=1, f(2n)=2f(n), f(2n+1)=2n+1+2f(n)+f(n)/n$. $S(n)=\sum f(i)^2$. Given $S(10)=1530, S(100)=4798445$. Find $S(10^{{16}}) \bmod 10^9+7$.

\section{Mathematical Analysis}
The recurrence suggests $f$ is related to a divide-and-conquer structure on binary representations. Since $f(2n)=2f(n)$, the function doubles with even arguments, similar to ruler functions.

For the sum of squares, we can decompose based on even/odd indices and use the binary structure to compute $S(N)$ in $O(\log N)$ recursive calls, each doing $O(1)$ work.

\section{Answer}

\[
\boxed{282771304}
\]

\end{document}
